\section{Limitaciones}
\label{sec:limitations}

Aunque el estudio presenta resultados positivos, existen limitaciones que merecen consideración. El dataset utilizado, pese a ser representativo, procede de una única fuente pública, circunstancia que puede comprometer la generalización hacia otros centros médicos equipados con instrumentos y protocolos de adquisición distintos. No se realizó validación externa en cohortes independientes ni se contrastó el rendimiento en conjuntos de datos multicéntricos. Además, todas las arquitecturas fueron entrenadas desde cero sin aprovechamiento de aprendizaje por transferencia, incrementando los tiempos de convergencia en relación con enfoques que utilizan modelos preentrenados. Finalmente, la investigación se restringió a la categorización de cuatro tipos tumorales; la detección temprana de neoplasias de bajo grado y la delineación precisa de márgenes tumorales escapan al alcance de este trabajo.
